% Drop-in replacement for the experiment section.
% Required packages: booktabs, tabularx, graphicx.
% The main-table numbers and Figures 1--5 are intentionally placeholders until
% the frozen p8 results are regenerated.  Existing q00_1096 figures average
% prefixes and therefore must not be relabeled as p8 results.

\providecommand{\TBD}{\textit{TBD}}
\providecommand{\model}[1]{\textsc{#1}}
\providecommand{\figplaceholder}[2][3.1cm]{%
  \fbox{\parbox[c][#1][c]{0.94\linewidth}{\centering\small #2}}}

\section{Experiments}
\label{sec:experiments}

We ask three questions: which spatial capabilities are accessible in a frozen
video representation, where in the network they emerge, and when additional
temporal evidence helps.  All probes operate on causal ScanNet++ prefixes and
use scene-disjoint splits.  Unless stated otherwise, the backbone is frozen,
the prefix contains eight observations (p8), and each model is evaluated at one
registered representative layer.  Longer prefixes and layer sweeps are analyzed
separately, rather than selected post hoc for the headline table.

\subsection{Main benchmark}
\label{sec:main-benchmark}

\paragraph{Capabilities and metrics.}
Table~\ref{tab:benchmark-tasks} summarizes the six tasks.  We retain the short
codes only as table indices; throughout the text we refer to the capabilities
by their scientific names.  The geometry tasks test metric scene structure and
cross-view correspondence.  Hidden-object persistence asks whether a queried
object remains localized after leaving the view.  The three dynamics tasks then
test one-step update, multi-step integration, and sensitivity to an alternative
camera action.

\begin{table*}[t]
  \centering
  \caption{\textbf{Benchmark axes and tasks.}  Targets are constructed from the
  same causal stream, while future target features used by the dynamics tasks
  are extracted from isolated target frames.  Higher is better except for depth
  loss and angular error.}
  \label{tab:benchmark-tasks}
  \small
  \begin{tabularx}{\textwidth}{@{}llXXXl@{}}
    \toprule
    ID & Capability & Probe input & Prediction target & What success shows & Metric \\
    \midrule
    A1 & Metric geometry & Streaming visual features & Depth and camera pose &
      Metric 3D structure is directly readable & Depth loss $\downarrow$; AUC@30 $\uparrow$ \\
    A2 & View correspondence & A pair of streaming views & View-overlap label &
      Scene identity survives viewpoint change & Balanced accuracy $\uparrow$ \\
    B2 & Hidden-object persistence & Object query and causal history &
      Object bearing in the final camera frame & An occluded object state is retained &
      Angular error (deg) $\downarrow$ \\
    C1 & One-step dynamics & Context and one 9D relative pose &
      Isolated next-state feature & The state can be updated by an action &
      Global R@1 $\uparrow$ \\
    C2 & Path integration & Context and a pose sequence &
      Features at successive waypoints & Egomotion can be accumulated &
      Global R@1 $\uparrow$ \\
    C3 & Counterfactual update & Shared context and alternative poses &
      Intervention-specific target features & Predictions change with the intervention &
      Intervention validity $\uparrow$ \\
    \bottomrule
  \end{tabularx}
\end{table*}

\paragraph{Models and default operating point.}
We compare ten frozen representations spanning image encoders, masked and
predictive video encoders, explicit 3D reconstruction, multimodal encoders, and
video generators: DINOv2-L, VideoMAE-L, V-JEPA2 ViT-L, Fast3R ViT-L,
Qwen2.5-VL-7B, Wan2.1-T2V-1.3B, CogVideoX-5B-I2V, AetherV1, MiniMax-H3, and
LingBot-Video MoE-30B-A3B.  Their representative layer indices are respectively
$0,23,23,24,-1,20,20,1,30,$ and $30$; indices $0$ for DINOv2 and $-1$ for
Qwen2.5-VL denote their single exported final representation and visual merger.
For diffusion models, the default extraction point is scheduler index 749 for
Wan, CogVideoX, and Aether; MiniMax-H3 and LingBot use flow coordinate $0.25$.

\begin{table*}[t]
  \centering
  \caption{\textbf{Main results at the fixed representative layer and p8.}
  Results are reported in native units; bold and underline indicate the best
  and second-best model per column.  No overall average is reported because the
  tasks measure distinct capabilities.  Cells marked \TBD{} await the frozen
  p8 export.}
  \label{tab:main-results}
  \setlength{\tabcolsep}{3.7pt}
  \small
  \begin{tabular}{@{}lrrrrrrr@{}}
    \toprule
    Model (layer) & Depth $\downarrow$ & Camera $\uparrow$ & View $\uparrow$ &
    Persistence $\downarrow$ & One-step $\uparrow$ & Integration $\uparrow$ &
    Counterfactual $\uparrow$ \\
    \midrule
    \model{DINOv2-L} (0)       & \TBD & \TBD & \TBD & \TBD & \TBD & \TBD & \TBD \\
    \model{Fast3R} (24)        & \TBD & \TBD & \TBD & \TBD & \TBD & \TBD & \TBD \\
    \model{VideoMAE-L} (23)    & \TBD & \TBD & \TBD & \TBD & \TBD & \TBD & \TBD \\
    \model{V-JEPA2-L} (23)     & \TBD & \TBD & \TBD & \TBD & \TBD & \TBD & \TBD \\
    \model{Qwen2.5-VL} ($-1$) & \TBD & \TBD & \TBD & \TBD & \TBD & \TBD & \TBD \\
    \model{Wan} (20)           & \TBD & \TBD & \TBD & \TBD & \TBD & \TBD & \TBD \\
    \model{CogVideoX} (20)     & \TBD & \TBD & \TBD & \TBD & \TBD & \TBD & \TBD \\
    \model{Aether} (1)         & \TBD & \TBD & \TBD & \TBD & \TBD & \TBD & \TBD \\
    \model{MiniMax-H3} (30)    & \TBD & \TBD & \TBD & \TBD & \TBD & \TBD & \TBD \\
    \model{LingBot-Video} (30) & \TBD & \TBD & \TBD & \TBD & \TBD & \TBD & \TBD \\
    \bottomrule
  \end{tabular}
\end{table*}

\paragraph{A spatial capability atlas.}
Raw metrics cannot be averaged because they have different units and
directions.  Figure~\ref{fig:capability-atlas} therefore converts each column of
Table~\ref{tab:main-results} to a within-task percentile after orienting all
metrics so that larger is better.  The resulting atlas reveals specialization:
geometry, persistence, and action-conditioned prediction do not induce a single
model ranking.  In particular, strong direct geometry need not imply strong
dynamics, and the best model changes across one-step, integration, and
counterfactual prediction.

\begin{figure*}[t]
  \centering
  \IfFileExists{figures/paper/fig01_capability_atlas_p8.pdf}{%
    \includegraphics[width=0.92\textwidth]{figures/paper/fig01_capability_atlas_p8.pdf}%
  }{\figplaceholder[5.0cm]{\textbf{Figure 1 placeholder: Spatial Capability Atlas}\\[2mm]
      Rows: 10 models. Columns: Depth, Camera, View, Persistence, One-step,
      Integration, Counterfactual.\\
      Cell color: within-capability percentile at the fixed layer and p8.}}
  \caption{\textbf{Spatial Capability Atlas.}  Task-wise normalization is used
  only for visualization; all statistical comparisons use native metrics.
  This figure answers \emph{which model exposes which spatial capability?}}
  \label{fig:capability-atlas}
\end{figure*}

\paragraph{Geometry is not dynamics.}
Figure~\ref{fig:geometry-dynamics} averages normalized depth, camera, and view
scores on the horizontal axis and C1--C3 on the vertical axis, placing each
model at one fixed operating point.  The models do not collapse onto a common
diagonal, and the
leaders on the two axes differ.  Spatial competence is therefore better
described as a profile than as a scalar ``spatial score.''  B2 is kept separate
because persistence combines geometric readout with memory.

\begin{figure}[t]
  \centering
  \IfFileExists{figures/paper/fig02_geometry_vs_dynamics_p8.pdf}{%
    \includegraphics[width=\columnwidth]{figures/paper/fig02_geometry_vs_dynamics_p8.pdf}%
  }{\figplaceholder[4.1cm]{\textbf{Figure 2 placeholder: Geometry vs. Dynamics}\\[2mm]
      One labeled point per model.\\
      $x$: mean normalized A1/A2; $y$: mean normalized C1--C3.}}
  \caption{\textbf{Geometry versus dynamics.}  Marker color may encode model
  family and marker size B2 performance.  The plot tests whether spatial
  capability is well represented by one common scalar.}
  \label{fig:geometry-dynamics}
\end{figure}

\subsection{Where and when do capabilities emerge?}
\label{sec:layer-prefix-analysis}

\paragraph{Layer-wise emergence.}
We next restore all registered layers and plot a separate curve for each prefix
in Figure~\ref{fig:layer-emergence}.  Layer depth is normalized within a model,
and scores are normalized within a model--task pair; neither the layer nor the
prefix is selected by validation performance.  The curves are task-specific
and frequently non-monotonic.  Across V-JEPA2, VideoMAE, and Wan, geometry tends
to peak deeper than one-step dynamics, which often favors early or middle
blocks.  This is a tendency rather than a universal ordering: counterfactual and
persistence readouts can peak later, and the maximizing layer can change with
the prefix.

\begin{figure*}[t]
  \centering
  \IfFileExists{figures/paper/fig03_layer_emergence_by_prefix.pdf}{%
    \includegraphics[width=0.98\textwidth]{figures/paper/fig03_layer_emergence_by_prefix.pdf}%
  }{\figplaceholder[5.0cm]{\textbf{Figure 3 placeholder: Capability emergence across layers}\\[2mm]
      Small multiples by capability; $x$: normalized layer depth; $y$: oriented,
      within-sweep score.\\
      Four lines per panel: p8, p12, p16, p24; one row per multi-layer model.}}
  \caption{\textbf{Capability emergence across layers.}  Multiple prefixes are
  shown explicitly rather than averaged.  The figure asks where a capability
  becomes readable and whether that location is stable under additional
  temporal evidence.}
  \label{fig:layer-emergence}
\end{figure*}

\paragraph{Longer history is capability-dependent.}
At each model's fixed representative layer, Figure~\ref{fig:prefix-integration}
reports the oriented difference between p24 and p8, standardized by the
task-level score variation.  Most model--task cells are non-positive: extra
frames commonly saturate or interfere with hidden-object localization and
action readout.  The effect is nevertheless structured rather than uniformly
negative.  At V-JEPA2's fixed layer, camera geometry and counterfactual update
provide positive cases, while depth and view do not improve from p8 to p24.
Their full layer--prefix surfaces nevertheless contain benefits at other
depths, foreshadowing the interaction below.  Thus, long history is useful only
when the representation and readout integrate the relevant temporal evidence.

\begin{figure*}[t]
  \centering
  \IfFileExists{figures/paper/fig04_prefix_delta_p24_minus_p8.pdf}{%
    \includegraphics[width=0.90\textwidth]{figures/paper/fig04_prefix_delta_p24_minus_p8.pdf}%
  }{\figplaceholder[4.8cm]{\textbf{Figure 4 placeholder: Prefix integration}\\[2mm]
      Model $\times$ capability heatmap at each model's fixed layer.\\
      Cell value: oriented $(\mathrm{p24}-\mathrm{p8})/$ task standard deviation;
      red hurts, blue helps.}}
  \caption{\textbf{Prefix integration.}  Positive values mean that adding
  observations improves the task.  Raw differences are standardized only to
  share a color scale.  This figure asks whether more history is reliably useful.}
  \label{fig:prefix-integration}
\end{figure*}

\paragraph{Layer--prefix interaction.}
The one-dimensional summaries hide cases in which temporal evidence changes
which layer is most useful.  Figure~\ref{fig:layer-prefix-cases} shows the full
layer--prefix surface for four prespecified case studies.  V-JEPA2 depth and
hidden-object persistence are especially informative because their best layer
changes across prefixes; Wan persistence illustrates a middle-layer optimum,
whereas CogVideoX persistence illustrates degradation with longer prefixes.
These surfaces show that representation depth and temporal evidence are not
independent experimental knobs.

\begin{figure*}[t]
  \centering
  \IfFileExists{figures/paper/fig05_layer_prefix_cases.pdf}{%
    \includegraphics[width=0.98\textwidth]{figures/paper/fig05_layer_prefix_cases.pdf}%
  }{\figplaceholder[5.2cm]{\textbf{Figure 5 placeholder: Layer $\times$ Prefix case studies}\\[2mm]
      Four heatmaps: V-JEPA2--Depth, V-JEPA2--Persistence,
      Wan--Persistence, CogVideoX--Persistence.\\
      $x$: p8/p12/p16/p24; $y$: layer; color: within-panel oriented score;
      mark the best layer in every column.}}
  \caption{\textbf{Layer--prefix interactions.}  Each panel is normalized
  independently and should be interpreted by its shape, not by comparing color
  values across tasks.}
  \label{fig:layer-prefix-cases}
\end{figure*}

\paragraph{Main findings.}
The main benchmark supports four conclusions.  First, spatial capabilities are
specialized, so there is no universal ``strong-model'' ordering.  Second,
readable spatial structure emerges at task-dependent and often non-monotonic
depths.  Third, longer context is not uniformly beneficial: prediction and
consistency sometimes improve, whereas localization and action readout often
saturate or degrade.  Finally, the effect of history depends on representation
depth; in V-JEPA2, for example, the optimal depth can move as the prefix grows.

\subsection{Analysis experiments}
\label{sec:analysis-experiments}

The following experiments test whether the main scores reflect the intended
mechanism rather than persistence, action priors, or feature-space shortcuts.

\subsubsection{Does the probe use history and action?}
\label{sec:analysis-controls}

\paragraph{Action-conditioned prediction beyond persistence.}
For Wan we fix p12, layer 20, and diffusion index 749.  A two-layer predictor
receives causal feature maps and the 9D relative camera pose and predicts the
spatially pooled feature of an isolated future frame.  Retrieval is performed
against the entire validation target pool.  Table~\ref{tab:wan-c1-controls}
shows that the correct action beats shuffled and zero-action evaluation, but
does not reliably beat copying the last observed state.  The evidence therefore
supports a weak action signal, not a strong dynamics claim.

\begin{table}[t]
  \centering
  \caption{\textbf{Wan one-step controls} (p12, layer 20, $t=749$; 276 samples,
  69 scenes).  Shuffle averages 100 cross-scene permutations.}
  \label{tab:wan-c1-controls}
  \small
  \begin{tabularx}{\columnwidth}{@{}Xrr@{}}
    \toprule
    Condition & Global R@1 (\%) & $\Delta$ vs. correct (pp) \\
    \midrule
    Correct action & 15.22 & -- \\
    Shuffled real action & 9.95 & $-5.26$ \\
    Evaluation-time zero action & 5.43 & $-9.78$ \\
    Copy last streaming state & 11.96 & $-3.26$ \\
    Copy prefix-pooled state & 1.09 & $-14.13$ \\
    \bottomrule
  \end{tabularx}
\end{table}

Across all 275 observed wrong actions per sample, the correct action improves
cosine similarity by $8.64\!\times\!10^{-4}$ with a scene-bootstrap 95\% CI of
$[6.25\!\times\!10^{-4},1.13\!\times\!10^{-3}]$.  By contrast, its R@1 gain
over last-state copy is $3.26$ points with CI $[-1.09,7.61]$ ($p=0.083$).
Correct $>$ shuffled alone is therefore insufficient to rule out temporal
smoothness.

\paragraph{What part of history supports persistence?}
We perform fixed-probe interventions on Wan B2 at the same layer and prefix.
Removing all but the current frame increases angular error by
$10.01^\circ$; removing complete frames in which the queried instance was
visible increases it by $6.59^\circ$; and retaining the query without any visual
history increases it by $13.55^\circ$ (Table~\ref{tab:b2-history-controls}).
History content therefore matters, including frames that previously contained
the target.  However, frame-order shuffling and reversal produce small changes
whose confidence intervals include zero, so these interventions establish
content dependence rather than a temporally ordered update rule.

\begin{table}[t]
  \centering
  \caption{\textbf{Wan hidden-object history controls.}  Error increases are
  paired against full history plus the correct query; confidence intervals use
  scene/trajectory-clustered bootstrap.}
  \label{tab:b2-history-controls}
  \small
  \begin{tabularx}{\columnwidth}{@{}Xrr@{}}
    \toprule
    Input intervention & $\Delta$ error & 95\% CI \\
    \midrule
    Current frame + query only & $+10.01^\circ$ & $[3.42,16.69]$ \\
    Remove target-visible frames & $+6.59^\circ$ & $[0.58,12.91]$ \\
    Query only; no visual history & $+13.55^\circ$ & $[7.31,19.61]$ \\
    \bottomrule
  \end{tabularx}
\end{table}

Query masking, same-scene query swapping, frame-order shuffling, and history
reversal complete the control battery.  These are inference-time diagnostics:
zeroing a query feature is not equivalent to retraining a probe without a query
token, and a shuffled sequence is not a physically observed trajectory.

\subsubsection{Spatial information across denoising time}
\label{sec:timestep-analysis}

Video generators expose representations at a chosen noise or flow coordinate.
We therefore repeat the benchmark for Wan and CogVideoX at scheduler indices
$\{499,999\}$, with the default 749 as reference, and for MiniMax-H3 at flow
coordinates $\{0.5,1.0\}$, with the default $0.25$ as reference.  The canonical
extraction grid contains layers $\{15,20,25\}$ and prefixes
$\{8,12,16,24\}$ for both streaming and isolated-target features.  The primary
comparison fixes layer 20 and p8; the larger grid is used only to check whether
a timestep effect is stable across depth and history.

\begin{figure*}[t]
  \centering
  \IfFileExists{figures/paper/fig06_timestep_analysis.pdf}{%
    \includegraphics[width=0.95\textwidth]{figures/paper/fig06_timestep_analysis.pdf}%
  }{\figplaceholder[4.6cm]{\textbf{Analysis figure placeholder: Denoising timestep}\\[2mm]
      Three panels (Wan, CogVideoX, MiniMax-H3).\\
      $x$: scheduler/noise coordinate; $y$: normalized capability score;
      lines: geometry, persistence, dynamics.}}
  \caption{\textbf{Spatial information across denoising time.}  This analysis
  tests whether spatial information is concentrated at a particular generation
  stage.  Feature extraction/probing is in progress; no timestep conclusion is
  drawn from the current artifacts.}
  \label{fig:timestep-analysis}
\end{figure*}

\subsubsection{Can a frozen representation support action output?}
\label{sec:action-output}

To close the perception--action loop, we construct a binary decision task from
the p8 stream: given a current state, a goal, and two candidate 9D actions,
select the action that makes progress toward the goal.  The current completed
  pilot uses independently extracted DINOv2-L features at its single exported
  layer, 512 training examples, and 128 validation examples.
Table~\ref{tab:action-output} separates endpoint recognition from action
selection and exposes action-vector shortcuts.

\begin{table}[t]
  \centering
  \caption{\textbf{DINOv2 p8 action-output pilot.}  All probes are trained from
  scratch on the same frozen split.}
  \label{tab:action-output}
  \small
  \begin{tabularx}{\columnwidth}{@{}Xr@{}}
    \toprule
    Available information & Accuracy (\%) \\
    \midrule
    Endpoint-feature oracle: goal + two realized endpoints & 69.5 \\
    Two candidate actions only & 64.1 \\
    Current context + anchor only & 56.3 \\
    Current + anchor + goal + two actions & 50.0 \\
    Random choice & 50.0 \\
    \bottomrule
  \end{tabularx}
\end{table}

The endpoint oracle shows that the representation can compare realized outcomes
to a goal.  Yet action-only prediction is already strong, while the direct
goal-conditioned actor remains at chance.  The pilot therefore reveals action
statistics and endpoint utility, but provides no evidence that the frozen
representation supports compositional transition reasoning.  A matched Wan
p8/p12 replication should be reported separately when available.

\subsubsection{Does V-JEPA2 propagate readable spatial structure?}
\label{sec:predicted-feature-probing}

The main benchmark probes encoder states and asks whether spatial information is
present.  For a predictive world model, a stronger requirement is that the
predictor carry the same information into the future.  We propose a
read-then-predict test.  First, a spatial readout $g$ is trained only on real
future encoder features and then frozen.  We apply this identical readout to the
real target, V-JEPA2's predicted future feature, a no-condition prediction, and
the copied last state.  The target-feature score is the readout ceiling; spatial
retention is credited only when the prediction improves over both controls.

\begin{figure*}[t]
  \centering
  \IfFileExists{figures/paper/fig07_vjepa_read_then_predict.pdf}{%
    \includegraphics[width=0.92\textwidth]{figures/paper/fig07_vjepa_read_then_predict.pdf}%
  }{\figplaceholder[4.6cm]{\textbf{Analysis figure placeholder: Read--Then--Predict}\\[2mm]
      Context $\rightarrow$ V-JEPA2 predictor $\rightarrow$ predicted future latent.\\
      Real target / predicted / no-condition / copy $\rightarrow$ one frozen
      depth--pointmap readout $\rightarrow$ spatial score.}}
  \caption{\textbf{Read--Then--Predict.}  The same frozen readout separates
  future-feature retrieval from preservation of spatially readable structure.}
  \label{fig:read-then-predict}
\end{figure*}

\begin{table}[t]
  \centering
  % Implementation audit: features/vjepa2/vjepa2_feature.py currently calls the
  % model with skip_predictor=True, so no result should be inserted here until a
  % predictor-feature export is completed under the same protocol.
  \caption{\textbf{Predicted-feature spatial retention.}  Results are pending
  predictor-feature extraction under the frozen protocol.}
  \label{tab:predicted-feature-retention}
  \small
  \begin{tabular}{@{}lrrr@{}}
    \toprule
    Input to frozen readout & Depth $\downarrow$ & Pointmap $\downarrow$ & vs. copy \\
    \midrule
    Real future feature & \TBD & \TBD & -- \\
    Predicted future feature & \TBD & \TBD & \TBD \\
    No-condition prediction & \TBD & \TBD & \TBD \\
    Last-state copy & \TBD & \TBD & -- \\
    \bottomrule
  \end{tabular}
\end{table}

\subsubsection{Does latent update preserve transformation order?}
\label{sec:se3-analysis}

Finally, we test a limited compositional property of action-conditioned latent
updates.  From a real observed triplet $A\!\rightarrow\!B\!\rightarrow\!C$,
we form continuous relative poses $T_1=T_{AB}$ and $T_2=T_{BC}$.  A predictor is
trained only on observed one-step pairs and then recursively applied as
$T_1$ followed by $T_2$.  At test time we compare its distance to the real
isolated $z_C$ with that obtained under the reversed order $T_2,T_1$.  This
tests whether latent updates retain action order; it is neither novel-view
rendering nor planning.

\begin{table*}[t]
  \centering
  \caption{\textbf{Exploratory observed-order pilot at p8.}  Cosine distance is
  averaged over seeds 42/43/44; lower is better.  The common held-out set has 44
  tuples from 11 scenes.}
  \label{tab:se3-pilot}
  \small
  \begin{tabular}{@{}lrrrrr@{}}
    \toprule
    Model (layer) & Correct order $\downarrow$ & Reversed $\downarrow$ &
    Margin $\uparrow$ & Order acc. $\uparrow$ & 95\% CI for margin \\
    \midrule
    DINOv2 (0)    & 0.2002 & 0.2040 & 0.0037 & 0.5455 & $[0.0004,0.0074]$ \\
    Fast3R (24)   & 0.0084 & 0.0087 & 0.0003 & 0.6288 & $[0.0000,0.0008]$ \\
    V-JEPA2 (23)  & 0.1704 & 0.1732 & 0.0028 & 0.4697 & $[-0.0028,0.0101]$ \\
    CogVideoX (20)& 0.0028 & 0.0029 & 0.0001 & 0.5152 & $[-0.0002,0.0005]$ \\
    \bottomrule
  \end{tabular}
\end{table*}

The margins are very small.  DINOv2 has the clearest positive interval; the
Fast3R margin is near numerical zero, while the V-JEPA2 and CogVideoX intervals
cross zero.  The pilot thus does not establish a general SE(3)-equivariant
latent representation.  It tests only two-step order sensitivity under the
motion families present in the data; rare motions, longer rollouts, and a
SeeSE3-style middle-layer replication remain missing.
